\scanpage{92}
% archive sheet number: 91

$\unsure{\mathbb V|V}$ \unsure{anneau de valuation discrète, complet}, à corps résiduel $k$ parfait, \unsure{corps des fractions} $K$, \unsure{uniformisante} $\pi$.

$\mathbb V$ le schéma en anneaux sur $k$ défini par $\mathbb V$

$\mathbb K$ le ind-schéma en anneaux défini par
\[
A\rightsquigarrow\mathbb V(A)_\pi=\mathbb K(A)
\]

\margin{$\unsure{\mathscr C_S|\mathscr E_S|C_S}$}
% slanted
\margin{\unsure{Notations}}
$S=\operatorname{Spec}(\unsure{V|\mathbb V})$, on munit $S$ de la ``topologie de Weil'' définie par
\[
\left\{\begin{aligned}
&\text{\unsure{catégorie des} $S'$ plats de type fini sur $S$.}\\
&\text{descente fid. plate [\uline{\unsure{ou fid. plt. finie} ???}]}
\end{aligned}\right.
\]
% slanted
\margin{\unsure{y-a-t-il lieu de prendre} $S'$ \unsure{simple} sur $S$? Non, \unsure{bien on va perdre les|bien on peut prendre les} \unsure{groupes|gr.} \unsure{infinitésimaux}, $S'$ \unsure{tel que} $S'_0$ \unsure{soit} \deleted{\unsure{réduit}} \unsure{réduit}?? Peut-être\unsure{\ldots}}

$C(S)=$ catégorie des faisceaux de Weil abéliens sur $S$

\unsure{Donner un} $G\in C(S)$ est un foncteur $S'\rightsquigarrow G(S')$, \unsure{soumis à des conditions d'exactitude connues}.

\unsure{On convient de} \uline{prolonger} $G$ \inserted{en $\overline G$} aux préschémas plats $S'$ quelconques sur $S$, par les conditions suivantes

\begin{enumerate}[label=(\roman*)]
\item $\overline G$ \unsure{de nature locale}
\item $\overline G$ \unsure{commute aux limites inductives d'anneaux}
\end{enumerate}

Alors \unsure{on aura} $\unsure{\overline m|=}$

\begin{enumerate}[label=(\roman*),start=3]
\item $\overline G$ \unsure{compatible avec descente par morphismes} fid. plats de présentation finie.
\end{enumerate}

Exemple : tout préschéma en gr. \inserted{\unsure{comm.} $G$} sur $S$ définit un faisceau de Weil, \unsure{dénoté par la} $=$ lettre. Si $G$ est loc. de type fini sur $S$, alors le \unsure{prolongement} can. $\overline G$ est encore ``représenté'' par le $\overline m$ $G$, i.e.
\[
\overline G(S')=\operatorname{Hom}_S(S',G).
\]
